%------------------------------------------------------------
% 3.2 Solución (Problema de dos cuerpos en CM + relativa)
%------------------------------------------------------------
\subsection*{3.2 \textcolor{Cinnabar}{\textbf{Solución:}}}

Partimos de las definiciones estándar para el problema de dos cuerpos:
\[
\vec{R}\;=\;\frac{m_1\vec{r}_1+m_2\vec{r}_2}{m_1+m_2},
\qquad
\vec{r}\;=\;\vec{r}_1-\vec{r}_2,
\qquad
M:=m_1+m_2.
\]

\paragraph{1) Expresar $\vec{r}_1$ y $\vec{r}_2$ en función de $\vec{R}$ y $\vec{r}$.}
De
\(
\vec{R}=\dfrac{m_1(\vec{r}_2+\vec{r})+m_2\vec{r}_2}{M}
\)
se obtiene, factorizando términos de $\vec{r}_2$,
\[
\vec{R}=\vec{r}_2+\frac{m_1}{M}\,\vec{r}
\;\;\Rightarrow\;\;
\boxed{\;\vec{r}_2=\vec{R}-\frac{m_1}{M}\,\vec{r}\;}
\]
y sustituyendo en
\(
\vec{r}=\vec{r}_1-\vec{r}_2
\)
queda
\[
\boxed{\;\vec{r}_1=\vec{R}+\frac{m_2}{M}\,\vec{r}\;}.
\]

Derivando en el tiempo,
\[
\dot{\vec{r}}_1=\dot{\vec{R}}+\frac{m_2}{M}\dot{\vec{r}},
\qquad
\dot{\vec{r}}_2=\dot{\vec{R}}-\frac{m_1}{M}\dot{\vec{r}}.
\]

\paragraph{2) Lagrangiana en coordenadas generalizadas $(\vec{R},\vec{r})$.}
La energía cinética total es
\[
T=\frac{1}{2}m_1\dot{\vec{r}}_1^{\,2}+\frac{1}{2}m_2\dot{\vec{r}}_2^{\,2}.
\]
Sustituyendo las expresiones anteriores y agrupando,
\[
T=\frac{1}{2}M\,\dot{\vec{R}}^{\,2}+\frac{1}{2}\,\mu\,\dot{\vec{r}}^{\,2},
\qquad
\boxed{\;\mu:=\frac{m_1m_2}{M}\;} \quad \text{(masa reducida)}.
\]
Para un potencial central $U(r)$ (y en particular $U=0$ en el caso libre) la
\textbf{Lagrangiana} queda
\[
\boxed{\;L=\frac{M}{2}\,\dot{\vec{R}}^{\,2}+\frac{\mu}{2}\,\dot{\vec{r}}^{\,2}-U(r)\;}
\]
y en el caso libre \(U=0\):
\[
\boxed{\;L=\frac{M}{2}\,\dot{\vec{R}}^{\,2}+\frac{\mu}{2}\,\dot{\vec{r}}^{\,2}\;}.
\]

\paragraph{3) Cantidades generalizadas y constantes de movimiento.}
Como $L$ no depende explícitamente de $\vec{R}$,
\[
\frac{d}{dt}\!\left(\frac{\partial L}{\partial\dot{\vec{R}}}\right)=\vec{0}
\;\;\Rightarrow\;\;
\boxed{\;\vec{P}:=\frac{\partial L}{\partial\dot{\vec{R}}}=M\dot{\vec{R}}=\text{cte.}\;}
\]
que es el \textit{momento lineal total}. Usando
\(
\dot{\vec{R}}=\dfrac{m_1\dot{\vec{r}}_1+m_2\dot{\vec{r}}_2}{M}
\),
\[
\boxed{\;\vec{P}=m_1\dot{\vec{r}}_1+m_2\dot{\vec{r}}_2=\vec{p}_1+\vec{p}_2\;}.
\]
En el sistema del centro de masa (\(\dot{\vec{R}}=\vec{0}\)) se tiene
\(\vec{P}=\vec{0}\Rightarrow m_1\dot{\vec{r}}_1+m_2\dot{\vec{r}}_2=\vec{0}\).

Para la coordenada relativa,
\[
\frac{\partial L}{\partial\dot{\vec{r}}}=\mu\dot{\vec{r}}=: \vec{p},
\qquad
\frac{d}{dt}\!\left(\frac{\partial L}{\partial\dot{\vec{r}}}\right)
-\frac{\partial L}{\partial\vec{r}}=\vec{0}
\;\;\Rightarrow\;\;
\mu\ddot{\vec{r}}+\nabla U(r)=\vec{0}.
\]
En el caso \textbf{libre} $U=0$, resulta \(\ddot{\vec{r}}=\vec{0}\) y por tanto
\(\dot{\vec{r}}=\) cte.\ y \(\vec{p}=\mu\dot{\vec{r}}=\) cte.

\paragraph{4) Conservación del momento angular orbital.}
Definimos el momento angular relativo
\[
\boxed{\;\vec{\ell}:=\vec{r}\times\vec{p}=\mu\,\vec{r}\times\dot{\vec{r}}\;}
\]
y su derivada temporal
\[
\dot{\vec{\ell}}
=\frac{d}{dt}\!\left(\vec{r}\times\mu\dot{\vec{r}}\right)
=\mu\big(\dot{\vec{r}}\times\dot{\vec{r}}+\vec{r}\times\ddot{\vec{r}}\big)
=\mu\,\vec{r}\times\ddot{\vec{r}}.
\]
Para fuerzas centrales \(\vec{F}=-\nabla U(r)=f(r)\,\hat{\vec{r}}\) se cumple
\(\vec{\tau}=\vec{r}\times\vec{F}=\vec{0}\), luego
\[
\boxed{\;\dot{\vec{\ell}}=\vec{0}\quad\Rightarrow\quad \vec{\ell}=\text{cte.}\;}
\]
y, en particular, en el caso libre (\(U=0\)), también \(\ddot{\vec{r}}=\vec{0}\Rightarrow\dot{\vec{\ell}}=\vec{0}\).

\paragraph{5) Resumen de resultados clave.}
\[
\boxed{\;\vec{r}_1=\vec{R}+\frac{m_2}{M}\vec{r},\qquad
       \vec{r}_2=\vec{R}-\frac{m_1}{M}\vec{r}\;}
\]
\[
\boxed{\;L=\frac{M}{2}\dot{\vec{R}}^{\,2}+\frac{\mu}{2}\dot{\vec{r}}^{\,2}-U(r),\quad
       M=m_1+m_2,\quad \mu=\frac{m_1m_2}{M}\;}
\]
\[
\boxed{\;\vec{P}=M\dot{\vec{R}}=\vec{p}_1+\vec{p}_2=\text{cte.},\qquad
       \vec{\ell}=\mu\,\vec{r}\times\dot{\vec{r}}=\text{cte.}\;}
\]
y en el sistema del centro de masa: \(\dot{\vec{R}}=\vec{0}\), \(\vec{P}=\vec{0}\).
